\section*{Chapitre \ref{ch:g11:enzymes-and-phenotype} --- Les enzymes et le phénotype}

\begin{solution}{exo:g11:enzymes-and-phenotype:1}
\omterm{def:g11:enzymes-and-phenotype:phenotype}{Génotype}~: l'ensemble des \omterm{def:g10:universal-dna:gene}{allèles} portés. \omterm{def:g11:enzymes-and-phenotype:phenotype}{Phénotype}~: l'ensemble des
caractéristiques observables, aux échelles moléculaire, cellulaire et
de l'organisme.
\end{solution}

\begin{solution}{exo:g11:enzymes-and-phenotype:2}
La poche de l'\omterm{def:g11:enzymes-and-phenotype:enzyme}{enzyme} où le \omterm{def:g11:enzymes-and-phenotype:enzyme}{substrat} se fixe et est transformé. La
spécificité~: le site s'ajuste à un \omterm{def:g11:enzymes-and-phenotype:enzyme}{substrat} (ou à une famille de
\omterm{def:g11:enzymes-and-phenotype:enzyme}{substrats}) et catalyse une seule réaction.
\end{solution}

\begin{solution}{exo:g11:enzymes-and-phenotype:3}
Vers \qty{37}{\celsius}~; à \qty{45}{\celsius}, il en reste environ un tiers.
\end{solution}

\begin{solution}{exo:g11:enzymes-and-phenotype:4}
Moléculaire~: tyrosinase absente ou inactive, aucune mélanine
fabriquée. Cellulaire~: \omterm{def:g10:cells-common-unit:cell}{cellules} pigmentaires présentes mais vides de
mélanine. Organisme~: cheveux blancs, peau et yeux très pâles,
sensibilité à la lumière et aux coups de soleil.
\end{solution}

\begin{solution}{exo:g11:enzymes-and-phenotype:5}
Le repliement de la pepsine ne fonctionne qu'en milieu très acide
(pH 2)~; l'intestin est neutre à légèrement alcalin, là où la courbe
montre une activité nulle.
\end{solution}

\begin{solution}{exo:g11:enzymes-and-phenotype:6}
Un seul \omterm{def:g10:universal-dna:gene}{allèle} fonctionnel produit la moitié de l'\omterm{def:g11:enzymes-and-phenotype:enzyme}{enzyme} normale, mais
une molécule d'\omterm{def:g11:enzymes-and-phenotype:enzyme}{enzyme} travaille des milliers de fois par seconde~: la
moitié de la quantité convertit encore assez de tyrosine pour faire un
pigment normal. Le \omterm{def:g11:enzymes-and-phenotype:phenotype}{phénotype} est saturé~; seuls deux \omterm{def:g10:universal-dna:gene}{allèles}
défaillants se voient.
\end{solution}

\begin{solution}{exo:g11:enzymes-and-phenotype:7}
Leur albinisme vient de blocages dans deux \omterm{def:g10:universal-dna:gene}{gènes} différents (deux
\omterm{def:g11:enzymes-and-phenotype:enzyme}{enzymes} différentes de la voie, ou du transport de la mélanine).
L'enfant hérite un \omterm{def:g10:universal-dna:gene}{allèle} fonctionnel de chaque \omterm{def:g10:universal-dna:gene}{gène} du parent dont le
blocage est ailleurs, si bien que les deux \omterm{def:g11:enzymes-and-phenotype:enzyme}{enzymes} fonctionnent et que
le pigment est fabriqué.
\end{solution}

\begin{solution}{exo:g11:enzymes-and-phenotype:8}
La phénylalanine s'accumule~; la tyrosine (et ce qui en dérive)
manque. Comme la mélanine est fabriquée à partir de la tyrosine, son
manque réduit le pigment~: cheveux et peau clairs sont fréquents dans
une phénylcétonurie non traitée.
\end{solution}

\begin{solution}{exo:g11:enzymes-and-phenotype:9}
Là où le lait est un aliment de base, le digérer apporte de l'énergie
et du calcium sans troubles, et l'\omterm{def:g10:universal-dna:gene}{allèle} de persistance est donc un
avantage. Là où les adultes ne boivent pas de lait, l'\omterm{def:g10:universal-dna:gene}{allèle} ne change
rien~; il n'est «~plus sain~» que dans l'environnement qui fournit le
\omterm{def:g11:enzymes-and-phenotype:enzyme}{substrat}.
\end{solution}

\begin{solution}{exo:g11:enzymes-and-phenotype:10}
Pour la tyrosinase~: moléculaire, le \omterm{def:g11:enzymes-and-phenotype:enzyme}{substrat} ne s'ajuste plus, aucune
mélanine~; cellulaire, des \omterm{def:g10:cells-common-unit:cell}{cellules} pigmentaires vides~; organisme,
l'albinisme. Pour la lactase~: le lactose n'est pas digéré~; les
\omterm{def:g10:cells-common-unit:cell}{cellules} intestinales sont privées du sucre et le côlon le fait
fermenter~; l'intolérance.
\end{solution}

\begin{solution}{exo:g11:enzymes-and-phenotype:11}
Ses extrémités restent au-dessus de \qty{35}{\celsius}, l'\omterm{def:g11:enzymes-and-phenotype:enzyme}{enzyme}
travaille à peine, et le chaton est pâle partout. Dehors en hiver, les
oreilles et les pattes se refroidissent, l'\omterm{def:g11:enzymes-and-phenotype:enzyme}{enzyme} devient active, et le
poil qui y repousse pousse foncé.
\end{solution}

\begin{solution}{exo:g11:enzymes-and-phenotype:12}
Moléculaire~: la moitié de l'hémoglobine est de la forme
drépanocytaire. Cellulaire~: à teneur normale en oxygène, les \omterm{def:g10:cells-common-unit:cell}{cellules}
restent rondes~; à très faible teneur, certaines se déforment en
faucille. Organisme~: en bonne santé dans la vie ordinaire, à risque en
altitude ou en plongée --- et protégé contre le paludisme.
L'environnement (l'oxygène, le paludisme) décide si le \omterm{def:g11:enzymes-and-phenotype:phenotype}{phénotype} du
porteur est un fardeau ou un avantage.
\end{solution}

\begin{solution}{exo:g11:enzymes-and-phenotype:13}
À \qty{41}{\celsius}, les \omterm{def:g11:enzymes-and-phenotype:enzyme}{enzymes} conservent l'essentiel de leur
activité (la courbe est près de son sommet)~; à \qty{43}{\celsius},
l'activité chute brutalement à mesure que les \omterm{def:g10:chemistry-of-life:families}{protéines} commencent à se
déplier, et si le dépliement se poursuit il est irréversible~: les
\omterm{def:g11:enzymes-and-phenotype:enzyme}{enzymes} sont détruites et non simplement ralenties.
\end{solution}

\begin{solution}{exo:g11:enzymes-and-phenotype:14}
Davantage de \omterm{def:g10:cells-common-unit:organelle}{mitochondries} et de capillaires par fibre musculaire~; un
cœur plus gros, au volume d'éjection plus grand~; plus d'hémoglobine
par litre de sang~; des réserves de \omterm{def:g10:exercise-and-energy:glycogen}{glycogène} plus grandes. Chacune est
un \omterm{def:g11:enzymes-and-phenotype:phenotype}{phénotype} produit par le même \omterm{def:g11:enzymes-and-phenotype:phenotype}{génotype} dans un environnement
différent~: les caractères dits «~génétiques~» sont la gamme que le
\omterm{def:g11:enzymes-and-phenotype:phenotype}{génotype} autorise, non une valeur fixée.
\end{solution}

\begin{solution}{exo:g11:enzymes-and-phenotype:15}
Le \omterm{def:g11:enzymes-and-phenotype:phenotype}{génotype} fixe quelle \omterm{def:g11:enzymes-and-phenotype:enzyme}{enzyme} est fabriquée et comment elle se
comporte~: pas d'\omterm{def:g11:enzymes-and-phenotype:enzyme}{enzyme} 1 (phénylcétonurie), pas de tyrosinase
(albinisme), une tyrosinase thermosensible (siamois). L'environnement
ne peut changer le résultat que dans les limites que l'\omterm{def:g11:enzymes-and-phenotype:enzyme}{enzyme} autorise~:
un régime retire le \omterm{def:g11:enzymes-and-phenotype:enzyme}{substrat} et sauve le cerveau, mais ne peut pas
fabriquer l'\omterm{def:g11:enzymes-and-phenotype:enzyme}{enzyme} manquante~; le froid fait travailler l'\omterm{def:g11:enzymes-and-phenotype:enzyme}{enzyme}
siamoise, mais aucune température ne fait travailler celle de
l'albinos~; et rien dans l'environnement ne donne son pigment à
l'albinos. Le \omterm{def:g11:enzymes-and-phenotype:phenotype}{génotype} fixe la gamme~; l'environnement choisit dedans.
\end{solution}

\begin{solution}{pb:g11:enzymes-and-phenotype:1}
\textbf{1.} Tronc \qty{38}{\celsius}~: environ 2~\%, pâle. Pattes
\qty{36}{\celsius}~: environ 20~\%, pâle. Extrémités des pattes
\qty{33}{\celsius}~: 80~\%, foncé. Oreilles et museau
\qty{31}{\celsius}~: plus de 90~\%, foncé.

\textbf{2.} Corps et pattes pâles, extrémités des pattes, oreilles,
museau (et queue) foncés~: le patron siamois.

\textbf{3.} 100~\% partout~: un chat uniformément pigmenté.

\textbf{4.} Dans l'utérus, chaque partie du chaton est à
\qty{38}{\celsius}, l'\omterm{def:g11:enzymes-and-phenotype:enzyme}{enzyme} est inactive partout, et aucun pigment
n'est déposé avant la naissance~; les extrémités foncent à mesure
qu'elles se refroidissent après la naissance.

\textbf{5.} Vers \qty{35.5}{\celsius}, là où l'activité franchit
30~\%~: c'est la partie descendante et raide de la courbe, entre 33 et
\qty{37}{\celsius}, qui fixe la limite.

\textbf{6.} Foncée~: à \qty{30}{\celsius}, l'\omterm{def:g11:enzymes-and-phenotype:enzyme}{enzyme} est pleinement
active dans les poils en croissance.

\textbf{7.} Pâle~: à \qty{38}{\celsius}, l'\omterm{def:g11:enzymes-and-phenotype:enzyme}{enzyme} est inactive.

\textbf{8.} Le pigment est déposé dans le poil à mesure qu'il pousse et
ne peut plus être ni retiré ni ajouté ensuite~; la couleur ne change
que lorsque les poils tombent et sont remplacés par de nouveaux poussés
à la température locale.

\textbf{9.} Le \omterm{def:g11:enzymes-and-phenotype:phenotype}{génotype} des \omterm{def:g10:cells-common-unit:cell}{cellules} de la plage est le même avant et
après, et le même que celui du reste du tronc~; seule la température de
l'\omterm{def:g11:enzymes-and-phenotype:enzyme}{enzyme} a changé, et le pigment avec elle. L'action porte sur
l'activité de l'\omterm{def:g11:enzymes-and-phenotype:enzyme}{enzyme}, un événement moléculaire.

\textbf{10.} Toutes les \omterm{def:g10:cells-common-unit:cell}{cellules} du chat descendent d'un même œuf et
portent les mêmes \omterm{def:g10:universal-dna:gene}{allèles}~; et les expériences de la plage rasée
changent la couleur d'une région sans changer ses \omterm{def:g10:cells-common-unit:cell}{cellules} --- une
poche de glace ne peut pas commuter un \omterm{def:g10:universal-dna:gene}{allèle}.

\textbf{11.} Moléculaire~: une tyrosinase qui ne travaille qu'en
dessous de \qty{35}{\celsius} environ. Cellulaire~: des \omterm{def:g10:cells-common-unit:cell}{cellules}
pigmentaires pleines de mélanine aux extrémités, vides sur le tronc.
Organisme~: un chat pâle à extrémités foncées.

\textbf{12.} L'\omterm{def:g10:chemistry-of-life:families}{acide aminé} occupe une position qui stabilise le
repliement~; le remplaçant tient le repliement à basse température mais
le laisse se relâcher quelques degrés plus haut, si bien que le \omterm{def:g11:enzymes-and-phenotype:enzyme}{site
actif} perd sa forme avant celui de l'\omterm{def:g11:enzymes-and-phenotype:enzyme}{enzyme} ordinaire.

\textbf{13.} L'\omterm{def:g11:enzymes-and-phenotype:enzyme}{enzyme} de l'\omterm{def:g10:universal-dna:gene}{allèle} ordinaire est pleinement active à la
température du corps~; la moitié de la quantité normale suffit encore
largement à fabriquer du pigment partout. L'\omterm{def:g10:universal-dna:gene}{allèle} siamois lui est
récessif.

\textbf{14.} Moléculaire~: aucune \omterm{def:g11:enzymes-and-phenotype:enzyme}{enzyme} fonctionnelle à quelque
température que ce soit, contre une \omterm{def:g11:enzymes-and-phenotype:enzyme}{enzyme} fonctionnelle sous
condition. Cellulaire~: aucune mélanine nulle part, contre de la
mélanine dans les régions fraîches. Organisme~: un chat entièrement
blanc aux yeux roses, contre le patron à extrémités foncées. Le froid
ne change rien pour l'albinos.

\textbf{15.} Qu'ils portent un \omterm{def:g10:universal-dna:gene}{allèle} thermosensible du même \omterm{def:g10:universal-dna:gene}{gène}~;
probablement une substitution voisine qui déstabilise l'\omterm{def:g11:enzymes-and-phenotype:enzyme}{enzyme},
apparue indépendamment dans chaque \omterm{def:g10:biodiversity-scales:species}{espèce}.

\textbf{16.} D'un facteur vingt~; l'\omterm{def:g11:enzymes-and-phenotype:enzyme}{enzyme} 1, qui convertit la
phénylalanine en tyrosine.

\textbf{17.} L'excès de phénylalanine gêne spécifiquement les \omterm{def:g10:cells-common-unit:cell}{cellules}
du cerveau en développement (leur approvisionnement en autres \omterm{def:g10:chemistry-of-life:families}{acides
aminés} et leur maturation)~; les \omterm{def:g10:cells-common-unit:cell}{cellules} de la peau le tolèrent. Le
défaut moléculaire est le même partout~; la conséquence cellulaire ne
l'est pas.

\textbf{18.} Le cinquième de \num{1200} fait
\qty{240}{\micro mol/L}, en dessous de \num{360}~: hors de danger.
Sans l'\omterm{def:g11:enzymes-and-phenotype:enzyme}{enzyme}, le taux sanguin suit simplement l'apport, de sorte que
le \omterm{def:g11:enzymes-and-phenotype:phenotype}{phénotype} est fixé par l'apport de \omterm{def:g11:enzymes-and-phenotype:enzyme}{substrat} --- l'environnement ---
et non par la seule \omterm{def:g11:enzymes-and-phenotype:enzyme}{enzyme} manquante.

\textbf{19.} La phénylalanine est nécessaire à la construction des
\omterm{def:g10:chemistry-of-life:families}{protéines} du corps et ne peut pas être synthétisée~; sans elle, la
croissance s'arrête. L'apport doit suffire à la synthèse des \omterm{def:g10:chemistry-of-life:families}{protéines}
et pas davantage, ce qui exige de mesurer et d'ajuster.

\textbf{20.} Le \omterm{def:g11:enzymes-and-phenotype:phenotype}{génotype} a fixé l'\omterm{def:g11:enzymes-and-phenotype:enzyme}{enzyme} elle-même
(thermosensible chez le chat, absente chez l'enfant)~; l'environnement
en a réglé les conditions
(température cutanée~; phénylalanine alimentaire)~; le \omterm{def:g11:enzymes-and-phenotype:phenotype}{phénotype} est le
résultat du \omterm{def:g11:enzymes-and-phenotype:phenotype}{génotype} exprimé dans un environnement donné.
\end{solution}
